%=======================================================================
%  WISSENSCHAFTLICHES POSTER  --  Beamer + beamerposter
%  Layout nachgebaut nach dem JKU-Vorlagenposter (A0 hoch, 3 Spalten)
%
%  KOMPILIEREN:   ./build.sh          -> poster.pdf + preview.png
%                 ./build.sh watch    -> baut bei jedem Speichern neu
%                 ./build.sh check    -> misst nur das Layout nach
%                 (oder schlicht:  pdflatex poster.tex)
%
%  WO SCHREIBE ICH WAS?
%    Zeile  85 ff. : STELLSCHRAUBEN -- Farben und alle Abstaende
%    Zeile 165     : \dummyfig  -- Platzhalterbild (spaeter durch
%                    \includegraphics ersetzen)
%    Zeile 191 ff. : Titel / Autoren / Institut / Fusszeile
%    Zeile 241 ff. : DER INHALT -- Spalte 1, 2, 3
%
%  Jeder Kasten wird mit  \begin{block}{Ueberschrift} ... \end{block}
%  gesetzt. Reihenfolge und Anzahl der Bloecke pro Spalte einfach
%  aendern -- sie stapeln sich automatisch untereinander.
%
%  DAS RASTER
%  \sidemargin (Rand aussen), \colsep (zwischen den Spalten), \topgap
%  und \botgap (zum Kopf-/Fussbalken) sowie \blockskip (zwischen zwei
%  Bloecken) sind feste mm-Werte; \colw ergibt sich daraus, damit die
%  Zeile immer genau so breit wie das Papier ist.  \blockinnersep ist
%  der Abstand vom Kastenrand zum Text -- er gilt fuer Titelbalken und
%  Koerper gleichermassen, damit beide dieselbe Fluchtlinie haben.
%  Achtung: die Spalten sind normale minipages, KEINE beamer-columns --
%  die wuerden zwischen den Spalten eigenes \hfil einfuegen, wodurch
%  \colsep nicht mehr stimmt.  Zwischen den Spaltenstuecken darf daher
%  auch keine Leerzeile stehen (die erzeugt \par und bricht die Zeile).
%
%  BUENDIGES LAYOUT
%  Die Abstaende sind fest: \blockskip zwischen zwei Bloecken, \topgap
%  unter dem Kopfbalken, \botgap ueber dem Fussbalken (Abschnitt
%  STELLSCHRAUBEN).  Dass alle drei Spalten unten auf gleicher Hoehe
%  enden, ist dagegen eine Frage der Inhaltsmenge -- LaTeX streckt
%  nichts von selbst.  Der Dummy-Inhalt ist darauf abgestimmt.
%  Sobald du eigenen Text einsetzt, verschiebt sich das.  Dann:
%     ./build.sh check
%  ausfuehren -- es meldet fuer jede Spalte, um wieviele Millimeter sie
%  zu lang oder zu kurz ist.  Laeuft eine Spalte so weit ueber, dass sie
%  im Fussbalken verschwindet, hilft  ./measure_free.sh  -- das misst an
%  einer Kopie ohne Fussbalken und zeigt die echte Unterkante.  Ausgleichen ueber die Hoehe der Bilder,
%  ein paar Zeilen Text mehr/weniger, oder indem ein Block in eine
%  andere Spalte umzieht.
%
%  ZWEI STOLPERFALLEN (schon geloest, aber gut zu wissen):
%   * \usepackage{exscale} muss drin bleiben, sonst bleiben \sum, \int
%     und grosse Klammern winzig, weil sie nicht mit der Postergroesse
%     mitskalieren.
%   * Formeln, die breiter als die Spalte sind, laufen stumm aus dem
%     Kasten heraus.  Dann mit  \begin{multline*} ... \\ ... \end{multline*}
%     umbrechen (Beispiel in Spalte 2, Block "Ergebnis A").
%     Kontrolle:  ./build.sh  meldet jede zu breite Zeile.
%=======================================================================
\documentclass[final,t]{beamer}

%-----------------------------------------------------------------------
% Papierformat.  orientation=portrait -> Hochformat,  size=a0 -> 841x1189mm
% scale  = globaler Schriftgroessen-Faktor.  Text zu klein/gross?  -> hier drehen!
%-----------------------------------------------------------------------
\usepackage[orientation=portrait,size=a0,scale=1.055]{beamerposter}

\usepackage[utf8]{inputenc}
\usepackage{lmodern}          % Vektorschriften (sonst pixelige Bitmapfonts)
\usepackage[T1]{fontenc}
\usepackage[english]{babel}      % <- fuer deutsches Poster: {ngerman}
\usepackage{amsmath,amssymb,amsfonts}
\usepackage{exscale}   % WICHTIG: laesst \sum, \int, Klammern mit der Postergroesse mitwachsen
\usepackage{tikz}
\usepackage{graphicx}
\usepackage{ragged2e}
\usepackage{array,booktabs}
\usepackage{bbold}
\usetikzlibrary{calc}

\usetheme{default}
\setbeamertemplate{navigation symbols}{}   % keine Navigationsleiste
% beamer legt sonst eigene Seitenraender an; die Raender macht hier
% ausschliesslich \sidemargin.
\setbeamersize{text margin left=0pt,text margin right=0pt}
\setbeamertemplate{caption}[numbered]

%=======================================================================
%  S T E L L S C H R A U B E N  (Farben & Masse)
%=======================================================================
% --- Farben (aus dem Vorlagenposter abgegriffen) ------------------------
\definecolor{jkuheader} {RGB}{ 60, 94,124}   % helles Blau: Kopfbalken
\definecolor{jkudark}   {RGB}{ 29, 56, 80}   % dunkles Navy: Block-Titelbalken
\definecolor{jkufoot}   {RGB}{ 29, 56, 80}   % Fussbalken
\definecolor{blockbg}   {RGB}{255,255,255}   % Bloecke: weiss
\definecolor{blockframe}{RGB}{176,184,196}   % Rahmen der Platzhalterbilder
\definecolor{pagebg}    {RGB}{225,231,240}   % Posterhintergrund: leicht blaeulich
\definecolor{logocolor} {RGB}{  0,  0,  0}   % JKU-Logo im Kopfbalken (schwarz)

% --- Masse -------------------------------------------------------------
% Rand aussen und Abstand zwischen den Spalten werden vorgegeben, die
% Spaltenbreite ergibt sich daraus -- so bleibt die Zeile immer exakt
% so breit wie das Papier, egal was man hier einstellt.
\newlength{\sidemargin}\setlength{\sidemargin}{15mm}  % Rand links/rechts
\newlength{\colsep}    \setlength{\colsep}    {9mm}   % Abstand zwischen den Spalten
\newlength{\colw}      \setlength{\colw}%
  {\dimexpr(\paperwidth-2\sidemargin-2\colsep)/3\relax}   % ergibt sich
% Diese drei Werte bestimmen das Raster.  \topgap und \botgap sind
% bewusst groesser als \blockskip, damit die Blockreihe als ein Feld
% zwischen Kopf- und Fussbalken steht.
\newlength{\blockskip} \setlength{\blockskip} {9mm}              % Abstand zwischen zwei Bloecken
\newlength{\topgap}    \setlength{\topgap}    {15mm}             % Luft Kopfbalken -> erster Block
\newlength{\botgap}    \setlength{\botgap}    {15mm}             % Luft letzter Block -> Fussbalken
\newlength{\headerh}   \setlength{\headerh}   {0.072\paperheight}% Hoehe Kopfbalken
\newlength{\footh}     \setlength{\footh}     {0.022\paperheight}% Hoehe Fussbalken

\setbeamercolor{background canvas}{bg=pagebg}

%=======================================================================
%  B L O C K - D E S I G N   (dunkler Titelbalken + heller Koerper)
%=======================================================================
\setbeamercolor{block title}{fg=white,bg=jkudark}
\setbeamercolor{block body} {fg=black,bg=blockbg}
\setbeamerfont{block title}{size=\large,series=\bfseries}

% Innenabstand.  Derselbe Wert wird im Titelbalken (leftskip/rightskip)
% und im Koerper (\fboxsep) benutzt -- dadurch stehen Ueberschrift und
% Fliesstext auf derselben Fluchtlinie.
\newlength{\blockinnersep} \setlength{\blockinnersep}{9mm}

\setbeamertemplate{block begin}{%
  \vskip0pt%
  % --- Titelbalken: volle Spaltenbreite ---
  \begin{beamercolorbox}[wd=\linewidth,leftskip=\blockinnersep,
                         rightskip=\blockinnersep,ht=2.6ex,dp=1.1ex]{block title}%
    \usebeamerfont{block title}\insertblocktitle%
  \end{beamercolorbox}%
  \nointerlineskip%
  % --- Koerper: Satzspiegel = \textwidth minus 2x Innenabstand.
  %     Mit \fboxsep=\blockinnersep wird die farbige Flaeche darum herum
  %     wieder genau \textwidth breit, also buendig mit dem Titelbalken.
  \setbox0=\vbox\bgroup
  \hsize=\dimexpr\linewidth-2\blockinnersep\relax
  \linewidth=\hsize \columnwidth=\hsize \textwidth=\hsize
  % kein \leavevmode: der Koerper muss im vertikalen Modus beginnen,
  % sonst kann ein Block nicht mit \wrappic anfangen.
  \vskip-0.5ex\justifying\ignorespaces%
}
\setbeamertemplate{block end}{%
  \par\vskip1.2ex\egroup
  \noindent
  \setlength{\fboxsep}{\blockinnersep}%
  \colorbox{blockbg}{\box0}%
  \vskip\blockskip%
}

% Aufzaehlungszeichen wie im Original: kleines gefuelltes Dreieck
\setbeamertemplate{itemize item}{\raisebox{0.15ex}{\scriptsize$\blacktriangleright$}}
\setbeamertemplate{itemize subitem}{\raisebox{0.15ex}{\tiny$\blacktriangleright$}}
\setbeamercolor{itemize item}{fg=jkudark}
\setbeamercolor{itemize subitem}{fg=jkudark}
% Aufzaehlungen enger setzen (spart viel Platz auf dem Poster)
\setlength{\leftmargini}{2.2em}
\newcommand{\tightlist}{\setlength{\itemsep}{0.35ex}\setlength{\parskip}{0pt}\setlength{\topsep}{0.4ex}}

%=======================================================================
%  H I L F S M A K R O S
%=======================================================================
% Platzhalter-Grafik:  \dummyfig{Breite}{Hoehe}{Beschriftung}
% -> spaeter einfach ersetzen durch  \includegraphics[width=...]{bild.pdf}
\newcommand{\dummyfig}[3]{%
  \begin{tikzpicture}
    \draw[fill=white,draw=blockframe,line width=1.0pt] (0,0) rectangle (#1,#2);
    \draw[blockframe,line width=0.7pt] (0,0) -- (#1,#2);
    \draw[blockframe,line width=0.7pt] (0,#2) -- (#1,0);
    \node[align=center,text=gray,font=\small] at (#1/2,#2/2)
         {\colorbox{white}{\ #3\ }};
  \end{tikzpicture}}

% Eine Posterspalte.  Bewusst eine normale minipage statt beamers
% columns-Umgebung: die fuegt zwischen den Spalten eigenes \hfil ein,
% wodurch der Spaltenabstand nicht mehr dem eingestellten Wert
% entspricht.  \textwidth wird lokal auf die Spaltenbreite gesetzt,
% damit Angaben wie 0.55\textwidth innerhalb eines Blocks stimmen.
\newenvironment{postercolumn}
  {\begin{minipage}[t]{\colw}%
     \setlength{\textwidth}{\colw}\setlength{\columnwidth}{\colw}}
  {\end{minipage}}

% Kleine Bildunterschrift
% Etwas Luft davor, sonst klebt die Unterschrift am Fliesstext, wenn die
% Textspalte neben dem Bild hoeher ist als das Bild selbst.
\newcommand{\figcap}[1]{\vspace{0.9ex}{\footnotesize\itshape #1\par}\vspace{0.3ex}}

% ---------------------------------------------------------------------
% Weissrand, der in der Bilddatei selbst steckt (gemessen, in bp).
% Reihenfolge:  links unten rechts oben.  \pic und \picrow schneiden ihn
% weg, dadurch werden die Grafiken bei gleichem Platzbedarf groesser.
% Neues Bild ohne Eintrag -> wird ungeschnitten eingebunden.
% ---------------------------------------------------------------------
\expandafter\def\csname trim@Synergy.pdf\endcsname{8 30 22 12}
\expandafter\def\csname trim@PES.pdf\endcsname{0 16 0 5}
\expandafter\def\csname trim@sene_qc.pdf\endcsname{11 9 16 7}
\expandafter\def\csname trim@Dan_qc.pdf\endcsname{0 1 0 0}
\expandafter\def\csname trim@FMO7_HEOM_sene.pdf\endcsname{99 27 120 12}
\expandafter\def\csname trim@FMO7_Lindblad_sene.pdf\endcsname{99 27 120 12}
\expandafter\def\csname trim@qpu_200fs_ohne_reset.pdf\endcsname{8 9 7 8}
\newcommand{\picttrim}[1]{%
  \ifcsname trim@#1\endcsname\csname trim@#1\endcsname\else 0 0 0 0\fi}

% Bild einbinden, Weissrand abgeschnitten.  Die Optionen werden vorher mit
% \edef ausgerechnet -- graphicx expandiert die trim-Werte sonst nicht.
\newcommand{\incpic}[2][\linewidth]{%
  \begingroup
  \edef\ip@opt{width=\the\dimexpr#1\relax,trim=\picttrim{#2},clip}%
  \expandafter\includegraphics\expandafter[\ip@opt]{pictures/#2}%
  \endgroup}

% Bild aus dem Ordner pictures/, auf die aktuelle Zeilenbreite skaliert.
%   \pic{dateiname-ohne-endung}          -> volle Breite
%   \pic[0.8]{dateiname-ohne-endung}     -> 80 % der Breite, zentriert
\newcommand{\pic}[2][1.0]{%
  \begin{center}\incpic[#1\linewidth]{#2}\end{center}}

% Kleine Blochkugel, die zeigt, wie das Hadamard-Gatter den Zustand vom
% Nordpol auf den Aequator dreht.  Reines TikZ -- keine Bilddatei noetig.
% \R steuert die Groesse.
\newcommand{\blochgate}{%
\begin{tikzpicture}[line width=0.9pt,scale=1]
  \def\R{1.55}
  % Kugel und Aequator
  \draw[black!35] (0,0) circle (\R);
  \draw[black!35] (0,0) ellipse ({\R} and {0.34*\R});
  % Achsen (duenn, im Hintergrund)
  \draw[->,black!45] (0,0) -- (0,{1.28*\R}) node[above,black,font=\small]{$|0\rangle$};
  \draw[->,black!45] (0,0) -- (0,{-1.20*\R}) node[below,black,font=\small]{$|1\rangle$};
  \draw[->,black!45] (0,0) -- ({1.18*\R},0) node[right,black,font=\small]{$y$};
  \draw[->,black!45] (0,0) -- ({-0.98*\R},{-0.33*\R})
        node[anchor=east,black,font=\small]{$x$};
  % Zustand vor dem Gatter: Nordpol
  \draw[->,line width=2.2pt] (0,0) -- (0,{\R});
  % Zustand nach dem Gatter: +x auf dem Aequator
  \draw[->,line width=2.2pt,jkudark] (0,0) -- ({-0.72*\R},{-0.24*\R});
  \node[jkudark,font=\small,anchor=north east] at ({-0.52*\R},{-0.50*\R}) {$|{+}\rangle$};
  % die Drehung selbst
  \draw[->,dashed,line width=1.1pt,jkudark]
        (0,{1.04*\R}) to[out=180,in=72] ({-0.78*\R},{-0.20*\R});
  \node[jkudark,font=\small] at ({-1.02*\R},{0.66*\R}) {$H$};
\end{tikzpicture}}

% Formel auf EINE Zeile zwingen.  Verkleinert nur, wenn sie zu breit ist --
% ist sie schmal genug, bleibt die normale Schriftgroesse.
\newcommand{\fitline}[1]{%
  \par\smallskip\noindent
  \resizebox{\ifdim\width>\linewidth\linewidth\else\width\fi}{!}{$\displaystyle #1$}%
  \par\smallskip}

% Formel-Stichwort:  \why{Kurzkommentar}  -- steht kleiner unter einer Formel
\newcommand{\why}[1]{{\small\color{black!72}#1\par}\smallskip}

% Bildzeile: Text links, Bild rechts, beide oben buendig.
%   \picrow{0.54}{ ...Text... }{0.42}{FMO.pdf}
% Die beiden Breiten sollten zusammen ~0.97 ergeben, damit zwischen Text und
% Bild keine auffaellige Luecke bleibt.  Wichtig: die Hoehen der beiden
% Spalten grob angleichen (Bildhoehe = Breite/Seitenverhaeltnis), sonst
% entsteht Totraum.  Alles, was NACH \picrow kommt, steht unter dem Bild.
\newcommand{\picrow}[4]{%
  \begin{minipage}[t]{#1\linewidth}#2\end{minipage}\hfill
  \begin{minipage}[t]{#3\linewidth}%
    \vspace{-6mm}%   <- hier den negativen Wert MIT Einheit (z. B. mm) VOR dem Bild
    \incpic{#4}%
  \end{minipage}\par}

%=======================================================================
%  K O P F   und   F U S S
%=======================================================================
% ---- 1) HIER TITEL / AUTOREN / INSTITUT EINTRAGEN ---------------------
\newcommand{\postertitleA}{An Enhanced Algorithm for Simulating}      % 1. Titelzeile (kleiner)
\newcommand{\postertitleB}{Non-Unitary Dynamics on Quantum Computers} % 2. Titelzeile (gross)
\newcommand{\posterauthors}{\underline{Simon Mader}, Akhil Bhartiya, Tobias Kramer}
\newcommand{\posterinstitute}{Institute for Theoretical Physics, Johannes Kepler University Linz}
% ---- Fusszeile --------------------------------------------------------
% <-- HIER deine Kontaktadresse eintragen (steht bewusst noch leer)
\newcommand{\posterfooter}{Institute for Theoretical Physics, JKU Linz -- \textit{Kontakt: simon.mader@jku.at}}

\setbeamertemplate{headline}{%
  \begin{beamercolorbox}[wd=\paperwidth]{empty}
  \begin{tikzpicture}[remember picture,overlay]
    % blauer Balken ueber die volle Breite
    \fill[jkuheader] (0,0) rectangle (\paperwidth,-\headerh);
    % ---- LOGO: linksbuendig im Balken -------------------------------
    % Farbe ueber \logocolor (unten bei den Farben) steuerbar.
    % Echtes Logo:  \node[anchor=west] at (\sidemargin,-0.5\headerh)
    %                    {\includegraphics[height=0.55\headerh]{jku_logo.pdf}};
    \node[anchor=west,align=left,text=logocolor] at (\sidemargin,-0.46\headerh)
      {\fontsize{80}{80}\selectfont\bfseries JKU\\[1.4ex]
       \fontsize{20}{22}\selectfont\bfseries JOHANNES KEPLER\\[0.4ex]
       \fontsize{20}{22}\selectfont\bfseries UNIVERSIT\"AT LINZ};
    % ---- TITEL / AUTOREN: rechtsbuendig ------------------------------
    \node[anchor=north east,align=right,text=white,inner sep=0pt]
         at (\paperwidth-\sidemargin,-0.10\headerh)
      {\fontsize{54}{58}\selectfont\bfseries\postertitleA\\[0.2ex]
       \fontsize{74}{78}\selectfont\bfseries\postertitleB\\[1.2ex]
       \fontsize{48}{52}\selectfont\posterauthors\\[1.6ex]
       \fontsize{34}{36}\selectfont\posterinstitute};
  \end{tikzpicture}
  \vskip\headerh
  \end{beamercolorbox}
  % Luft zwischen Kopfbalken und erster Blockzeile.  Der Abzug faengt
  % den Innenabstand ab, den beamer selbst noch oberhalb des Rahmens
  % einfuegt -- so beginnt die erste Blockzeile exakt bei \topgap.
  \vskip\dimexpr\topgap-12.4mm\relax
}

\setbeamertemplate{footline}{%
  \begin{tikzpicture}[remember picture,overlay]
    \fill[jkufoot] (0,0) rectangle (\paperwidth,-\footh);
    \node[anchor=west,text=white,font=\fontsize{28}{30}\selectfont\bfseries]
         at (\sidemargin,-0.52\footh) {\posterfooter};
  \end{tikzpicture}
  \vskip\footh
}

\begin{document}
\begin{frame}[t]{}
% (Feinjustage der Startlinie: siehe \topgap oben)
\noindent\hspace*{\sidemargin}%
%=======================================================================
%%%%%%%%%%%%%%%%%%%%%%%%  S P A L T E   1  %%%%%%%%%%%%%%%%%%%%%%%%%%%%%
%  Warum ein Quantencomputer -> welches System -> welches Modell -> Qubits
%=======================================================================
\begin{postercolumn}

  %--- Bild: Synergy.pdf (Abb. 1.1) ------------------------------------
  \begin{block}{Motivation: Quantum and Classical Together}
    Open quantum systems exchange energy and coherence with their
    environment: their evolution is \textbf{non-unitary}. Quantum hardware
    executes only \textbf{unitary} gates -- this work bridges that gap.

    \pic[0.76]{Synergy.pdf}
    \figcap{Hybrid loop: the CPU builds and controls the circuit, the QPU
      carries the exponentially large state, the measurement is read back.}
    \smallskip

    \begin{itemize}\tightlist
      \item \textbf{Scaling:} $2^n$ amplitudes for $n$ qubits
      \item \textbf{Task:} embed a non-unitary map into a larger unitary one
      \item \textbf{Tool:} Kraus operators $+$ SVD $+$ dilation
      \item \textbf{Test system:} the FMO light-harvesting complex
    \end{itemize}
  \end{block}

  %--- Bild: Bloch_sphere.pdf (Abb. 5.1) -------------------------------
  \begin{block}{Qubits and the Bloch Sphere}
    \picrow{0.58}{%
      A qubit may be in a \textbf{superposition} of the basis states:
      \[ |\psi\rangle = \alpha|0\rangle + \beta|1\rangle \]
      \why{Normalised, $|\alpha|^2+|\beta|^2=1$, and $\Pr(0)=|\alpha|^2$.}
      With $\alpha=\cos\frac{\theta}{2}$, $\beta=e^{i\varphi}\sin\frac{\theta}{2}$
      the normalisation becomes two angles:
      \[ |\psi\rangle = \cos\tfrac{\theta}{2}|0\rangle
                      + e^{i\varphi}\sin\tfrac{\theta}{2}|1\rangle \]
    }{0.39}{Bloch_sphere.pdf}

    \textbf{A gate is a rotation of that arrow.} Probability conservation
    demands $U^\dagger U=\mathbb{1}$.
    Example -- the Hadamard gate applied to $|0\rangle$:
    \fitline{H = \tfrac{1}{\sqrt2}\begin{pmatrix}1&1\\1&-1\end{pmatrix},
             \quad
             H|0\rangle = \tfrac{1}{\sqrt2}\big(|0\rangle+|1\rangle\big)
                        \equiv |{+}\rangle}
    \begin{minipage}[t]{0.62\linewidth}
      \vspace{0pt}\why{$\theta:0\to\tfrac{\pi}{2}$ at $\varphi=0$: the arrow swings from
           the north pole onto the $+x$ axis -- a definite $|0\rangle$ turns
           into an equal superposition. Applying $H$ again brings it back,
           since $H|{+}\rangle=|0\rangle$.}
    \end{minipage}\hfill
    % Kugel bewusst hochgezogen und rechtsbuendig -- so bleibt unten Platz
    \begin{minipage}[t]{0.35\linewidth}
      \vspace{-25mm}\raggedleft\blochgate
    \end{minipage}
  \end{block}

  %--- Bild: FMO.pdf (Abb. 2.1) ----------------------------------------
  \begin{block}{The Fenna--Matthews--Olson Complex}
    \picrow{0.56}{%
      Seven bacteriochlorophyll (BChl) molecules guide an excitation from
      the antenna to the reaction centre~[1][3]. Absorption happens at
      $\approx 800\,$nm, in the infrared.
      \begin{itemize}\tightlist
        \item BChl 1, 6 at the antenna; BChl 3, 4 at the reaction centre
        \item the excitation \textbf{delocalises} over all seven sites,
              mediated by a dipole--dipole interaction -- no electron is
              ever exchanged
      \end{itemize}
    }{0.41}{FMO.pdf}
    \bigskip

    Each BChl is a \textbf{two-level system} (HOMO/LUMO). With at most one
    excitation in the complex this gives the site basis:
    \[ |i\rangle = |\phi^g_1\rangle\cdots|\phi^e_i\rangle\cdots|\phi^g_N\rangle \]
  \end{block}

  %--- Bild: PES.pdf (Abb. 2.2) -- Frenkel-Exziton-Modell ---------------
  \begin{block}{The Frenkel Exciton Model}
    Following Kreisbeck~[2] the Hamiltonian splits into three parts:
    \[ \hat H = \hat H_S + \hat H_B + \hat H_{SB} \]

    \textbf{System} -- site energies and dipole--dipole coupling:
    \[ \hat H_S = \sum_m \epsilon^e_m |m\rangle\langle m|
                + \sum_{m\neq n} J_{mn}\,|m\rangle\langle n| \]
    \why{$J_{mn}$: excitation hops $m\!\to\!n$ $\Rightarrow$ delocalisation.
         No electrons are exchanged.}

    \textbf{Bath} -- nuclear vibrations as harmonic normal modes:
    \[ \mathcal{H}_B = \sum_\xi \left( \frac{p_\xi^2}{2m_\xi} + 
    \frac{1}{2}m_\xi \omega_\xi^2 q_\xi^2 \right) = \sum_\xi \hbar\omega_\xi \left( b_\xi^\dagger b_\xi + 
    \frac{1}{2} \right), \]
    \why{$b^\dagger_\xi$ creates a phonon $\hbar\omega_\xi$ in mode $q_\xi$.}

    \begin{minipage}[t]{0.55\linewidth}
      \textbf{Coupling.} The nuclear dynamics follow potential energy
      surfaces (PES), harmonic around the equilibrium. In the excited state
      that equilibrium is shifted by $q_{\xi 0}$.
      \smallskip

      Expanding the shifted bath Hamiltonian leaves a term linear in
      $b^\dagger_\xi + b_\xi$ -- that is the coupling, and it exists only
      while the site is excited. Relaxation into the new minimum dissipates
      the \textbf{reorganisation energy}
      $\lambda_m=\sum_\xi \hbar\omega_{\xi m}d^2_{\xi m}$.
    \end{minipage}\hfill
    \begin{minipage}[t]{0.42\linewidth}
      \vspace{0pt}\incpic{PES.pdf}
      \figcap{Red arrow: Vertical Franck–Condon transition: nuclei remain fixed during excitation 
      ($\tau_\text{el} \ll \tau_\text{nuc}$), since electronic dynamics is much faster than nuclear motion.}
    \end{minipage}\par

    \bigskip
    \bigskip
    \bigskip

    \[ \hat H_{SB} = \sum_{m}\sum_\xi \hbar\omega_{\xi m}\, d_{\xi m}
          \big( b^\dagger_{\xi m} + b_{\xi m}\big)\otimes|m\rangle\langle m| \]
    \why{Displacement $d_\xi=-\sqrt{m_\xi\omega_\xi/2\hbar}\,q_{\xi 0}$;
         every mode couples to every site.}

  \end{block}%
\end{postercolumn}\hspace{\colsep}%
%=======================================================================
%%%%%%%%%%%%%%%%%%%%%%%%  S P A L T E   2  %%%%%%%%%%%%%%%%%%%%%%%%%%%%%
%  Markovsche Dynamik: Lindblad -> Kraus -> Schaltkreis -> Ergebnis
%=======================================================================
\begin{postercolumn}

  %---------------------------------------------------------------
  \begin{block}{Markovian Dynamics: Lindblad $\rightarrow$ Kraus}
    Lindblad equation - the most general Markovian master equation. Uses Born, Markov and secular approximation:
    \[ \frac{d}{dt}\rho_S = \underbrace{-\tfrac{i}{\hbar}[H,\rho_S]}_{\text{unitary}}
       + \underbrace{\sum_k \gamma_k \Big( L_k \rho_S L^\dagger_k
         - \tfrac12 \{ L^\dagger_k L_k, \rho_S \} \Big)}_{\text{dissipator} \rightarrow \text{non-unitary}} := \mathcal{L} \rho_S \]
    \why{$\gamma_k$: transition rate $\rightarrow$ Fermi's golden rule.
    $L_k$: Lindblad operator $\rightarrow$ jump operator.}

    \textbf{Kraus operators directly from it.} For a small step $dt$ the
    Ansatz
    \[ M_0 = \mathbb{1} + \Big(-\tfrac{i}{\hbar}H
             - \tfrac12\sum_k \gamma_k L^\dagger_k L_k\Big) dt,
       \qquad M_k = \sqrt{\gamma_k\,dt}\;L_k \]
    reproduces the Lindblad equation to $\mathcal{O}(dt)$
    \why{Hence it forces small steps. The algorithm therefore takes the route below, 
    which is exact for
         \emph{any} $t$.}

    \textbf{1. Vectorise} $\Rightarrow$ linear, integrable in closed form:
    \[ \mathrm{vec}(\rho_S(t)) = \Lambda(t)\,\mathrm{vec}(\rho_S(0)) \]
    \why{$\mathcal{L}$ is time independent $\Rightarrow$ the step $t$ may be
         \textbf{arbitrarily large}: no Trotterisation, no small-step error.}

    \textbf{2. Choi matrix} of the map, then diagonalise:
    \[ C(t) = \sum_{i,j}(E_{ij}\otimes \mathbb{1})\,\Lambda(t)\,(\mathbb{1} \otimes E_{ij})
            = \sum_k \sigma_k\, u_k u^\dagger_k \]
    \why{$\Lambda(t)$ is CPTP $\Rightarrow$ $C(t)$ positive, all
         $\sigma_k\ge 0$.}

    \textbf{3. Kraus operators} read off the eigenvectors:
    \[ \mathrm{vec}(M_k(t)) = \sqrt{\sigma_k(t)}\; u_k(t),
       \qquad \rho_S(t) = \sum_k M_k \rho_S(0) M^\dagger_k \]
  \end{block}

  %--- Bilder: sene_qc.pdf (Abb. 6.4), sene_Algo.pdf (Abb. 6.2) --------
  \begin{block}{Why the Circuit Looks the Way It Does}
    The $M_k$ are \textbf{not unitary}. An SVD isolates the non-unitary
    part in the diagonal matrix $\Sigma$:
    \[ M_k(t) = U\,\Sigma\,V^\dagger \]
    \why{$U$, $V^\dagger$ are rotations and already unitary -- only the
         rescaling $\Sigma$ must be repaired.}

    \textbf{Dilation} into a unitary of twice the size -- the price is
    exactly one ancilla qubit:
    \[ U_\Sigma = \begin{pmatrix} \Sigma_+ & 0\\ 0 & \Sigma_-\end{pmatrix},
       \qquad
       \Sigma_{\pm jj} = \sigma_{jj} \pm i\sqrt{1-\sigma_{jj}^2} \]
    \why{$|\sigma_{jj}|\le 1$ is guaranteed by complete positivity, so the
         root stays real.}

    \pic{sene_qc.pdf}
    \figcap{Seven-site FMO: three system qubits carry $|\psi\rangle$, the
      fourth is the ancilla, framed by two Hadamards.}

    Propagating $|\psi\rangle\otimes|0\rangle$ through the circuit and using
    $\Sigma_++\Sigma_-=2\Sigma$:
    \fitline{U \, H\,U_\Sigma\,V^\dagger H\,|\psi\rangle\otimes|0\rangle
      = \frac12\begin{pmatrix} U(\Sigma_++\Sigma_-)V^\dagger|\psi\rangle\\[0.4ex]
                               U(\Sigma_+-\Sigma_-)V^\dagger|\psi\rangle
               \end{pmatrix}
      = \begin{pmatrix} M_k(t)\,|\psi\rangle\\[0.4ex] |\varphi\rangle\end{pmatrix}}
    \why{\textbf{The upper half of the output is exactly $M_k(t)|\psi\rangle$};
         $|\varphi\rangle$ is a residue of the dilation and is discarded.}

    \pic{sene_Algo.pdf}
    \figcap{The loop: Lindblad $\rightarrow$ Choi $\rightarrow$ Kraus
      $\rightarrow$ SVD and dilation $\rightarrow$ circuit $\rightarrow$
      $\rho_S(t)$, then one step further in time.}
    \smallskip

    \textbf{The enhanced algorithm.} In Seneviratne \emph{et al.}~[4] the state is
    re-prepared classically at every time step. Here, the state of the
    \emph{previous} step is reused:
    \begin{itemize}\tightlist
      \item the \textbf{same} circuit is applied again and again
      \item the classical workload drops, the work moves to the QPU
      \item the growth of the classical cost with $t$ disappears, but QST is needed at every step
    \end{itemize}
  \end{block}

  %--- Bild: FMO7_Lindblad_sene.pdf (Abb. 6.5/6.7) ---------------------
  \begin{block}{Result: Markovian Dynamics}
    \pic[0.62]{FMO7_Lindblad_sene.pdf}
    \figcap{Site populations of the seven-site FMO complex over
      $1000\,$fs. Lines: classical reference, dots: quantum simulation (on a classical computer) 
      using qiskit. The population of site $i$ is the probability for the excitation to be on site $i$.}
  \end{block}%
\end{postercolumn}\hspace{\colsep}%
%=======================================================================
%%%%%%%%%%%%%%%%%%%%%%%%  S P A L T E   3  %%%%%%%%%%%%%%%%%%%%%%%%%%%%%
%  Nicht-Markovsche Dynamik: HEOM -> Schaltkreis + Ergebnis -> Hardware
%=======================================================================
\begin{postercolumn}

  %---------------------------------------------------------------
  \begin{block}{Non-Markovian Dynamics: HEOM}
    The Markov approximation assumes a bath without memory. The
    \textbf{hierarchical equations of motion} do not: the bath degrees of
    freedom are absorbed into \textbf{auxiliary density operators} (ADOs)
    $\hat\sigma_{\mathbf n}$ forming a coupled hierarchy,
    \smallskip
    \fitline{\frac{d}{dt}\hat\sigma_{\mathbf n}
       = \underbrace{-\frac{i}{\hbar}\mathcal{L}_S\hat\sigma_{\mathbf n}
         + \sum_{j,k}\big(\phi_j\theta_{j,0}-n_{jk}\gamma_{jk}\big)\hat\sigma_{\mathbf n}}
         _{\text{own dynamics and decay}}
       + \underbrace{\sum_{j,k}\phi_j\,\hat\sigma_{\mathbf n^+_{jk}}}_{\text{level above}}
       + \underbrace{\sum_{j,k} n_{jk}\,\theta_{j,k}\,\hat\sigma_{\mathbf n^-_{jk}}}_{\text{level below}}}
    \why{$\phi_j,\theta_{j,k}$: bath-coupling superoperators, $\gamma_{jk}$:
         decay rate of mode $jk$. Level $0$ is $\rho_S$ itself.}

    \textbf{Numerically exact} -- no Born, Markov or secular approximation;
    the only error is the truncation depth $\mathcal{N}$:
    \[ \#\,\mathrm{ADOs} = \frac{(\mathcal{N}+NK)!}{\mathcal{N}!\,(NK)!} \]
    \why{$N$ sites, $K$ bath exponentials -- grows fast, HEOM is expensive.}
  \end{block}

  %--- Bilder: Dan_qc.pdf (Abb. 6.18), FMO7_HEOM_sene.pdf (Abb. 6.20) --
  \begin{block}{The Non-Markovian Propagator on a Circuit}
    Following Dan \emph{et al.}~[5]: HEOM yields a propagator $G(t)$ whose
    $m$-th column holds the populations for an initial excitation on site
    $m$:
    \fitline{G(t) = \left\{
      \begin{pmatrix}
        \rho_{11}^1(t_0) & \rho_{11}^2(t_0) & \rho_{11}^3(t_0) & \rho_{11}^4(t_0)\\
        \rho_{22}^1(t_0) & \rho_{22}^2(t_0) & \rho_{22}^3(t_0) & \rho_{22}^4(t_0)\\
        \rho_{33}^1(t_0) & \rho_{33}^2(t_0) & \rho_{33}^3(t_0) & \rho_{33}^4(t_0)\\
        \rho_{44}^1(t_0) & \rho_{44}^2(t_0) & \rho_{44}^3(t_0) & \rho_{44}^4(t_0)
      \end{pmatrix},
      \begin{pmatrix}
        \rho_{11}^1(t_1) & \rho_{11}^2(t_1) & \rho_{11}^3(t_1) & \rho_{11}^4(t_1)\\
        \rho_{22}^1(t_1) & \rho_{22}^2(t_1) & \rho_{22}^3(t_1) & \rho_{22}^4(t_1)\\
        \rho_{33}^1(t_1) & \rho_{33}^2(t_1) & \rho_{33}^3(t_1) & \rho_{33}^4(t_1)\\
        \rho_{44}^1(t_1) & \rho_{44}^2(t_1) & \rho_{44}^3(t_1) & \rho_{44}^4(t_1)
      \end{pmatrix},\ \ldots\right\}}
    \why{Four-site case. Superscript: site of the initial excitation, hence
         $G(t_0)=\mathbb{1}$; $G$ is a \emph{list}, one matrix per time step.}
        \[ |\psi(t)\rangle = G(t)\,|\psi(0)\rangle, \quad G(t) = U\,\Sigma\,V^\dagger \]
    \why{$G(t)$ is non-unitary $\Rightarrow$ the same SVD and the same
         dilation $U_\Sigma$ as above.}

    \pic[0.51]{Dan_qc.pdf}
    \figcap{$U$, $V^\dagger$ come straight from the SVD of $G(t)$; $G(t)$
      itself is computed with HEOM.}

    \[U \, H\,U_\Sigma V^\dagger H\,|\psi(0)\rangle\otimes|0\rangle
       = \begin{pmatrix} |\psi(t)\rangle \\[0.4ex] |\varphi\rangle \end{pmatrix},
       \quad \rho_S = |\psi(t)\rangle\langle\psi(t)| \]

    \pic[0.60]{FMO7_HEOM_sene.pdf}
    \figcap{Seven-site FMO, non-Markovian: the coherent oscillation between
      sites 1 and 2 survives far longer than in the Markovian case.}
  \end{block}

  %--- Bild: qpu_200fs_ohne_reset.pdf -- durchgaengiges Gitter ---------
  \begin{block}{One Continuous Grid on Real Hardware}
    \picrow{0.67}{%
      \textbf{This work.} Stack $\rho_S$ and all ADOs into one super-vector
      $\vec\Sigma=(\mathrm{vec}\,\rho_S,\,\mathrm{vec}\,\sigma_1,\ldots)^{\!\top}$.
      There the evolution is exact and local in time -- the ADOs carry the
      memory, so the propagator needs no memory kernel and $P^n$ stays
      exact for every $n$.
      \smallskip

      \textbf{Obstacle.} A Sz.-Nagy dilation of $P$ costs a success
      probability $s^{-2n}$ with $s=\|P\|_2=32.39$: for 200 steps
      $10^{-301}$ -- hopeless, though the dynamics itself stays bounded.
    }{0.30}{qpu_200fs_ohne_reset.pdf}
    \figcap{Four-site FMO on \textsf{ibm\_kingston}, 70 circuits,
      $\Delta t=20\,$fs. Top: populations, bottom: coherence $\rho_{12}$;
      lines classical, squares hardware.}
    \[ \frac{d}{dt}\vec\Sigma = \mathcal{L}_{\mathrm{HEOM}}\vec\Sigma
       \;\Longrightarrow\;
       P = e^{\mathcal{L}_{\mathrm{HEOM}}\Delta t},
       \;\; \vec\Sigma_n = P^{\,n}\vec\Sigma_0 \]
    \why{$P$ forms a \textbf{semigroup}: $P^n$ is exact for every $n$.}

    \textbf{Fix: a metric in which $\mathcal{L}$ is dissipative.} $\|P\|_2=32$
    says something about the Euclidean inner product, not about the physics
    -- a Lyapunov solve gives a positive definite $W$:
    \[ \tilde P = W^{1/2} P\, W^{-1/2},
       \qquad \|\tilde P\|_2 \le e^{\delta\Delta t}\approx 1 \]
    \why{An \textbf{exact} similarity transform; success probability
         $\mathcal{O}(1)$.}

    \textbf{Compression is now free:} a Krylov projection of a contraction is
    again one, $\tilde H_m=\tilde Q^\dagger_m\tilde P\tilde Q_m$.
    \[ U = \begin{pmatrix}
             \tilde H_m/s & \sqrt{I-\tilde H_m \tilde H^\dagger_m/s^2}\\[0.4ex]
             \sqrt{I-\tilde H^\dagger_m \tilde H_m/s^2} & -(\tilde H_m/s)^\dagger
           \end{pmatrix} \]
    \why{$m=4$ $\Rightarrow$ \textbf{3 qubits}, 31 CNOTs per step:
         \textbf{one single unitary}, re-applied $n$ times. The classical
         cost is paid once, independent of the number of steps.}
  \end{block}

  %---------------------------------------------------------------
  \begin{block}{References}
    \scriptsize
    \begin{itemize}\tightlist
      \item[{[1]}] J.~Adolphs, T.~Renger, \emph{How Proteins Trigger
            Excitation Energy Transfer in the FMO Complex of Green Sulfur
            Bacteria}, Biophys.~J. \textbf{91}, 2778 (2006).
      \item[{[2]}] C.~Kreisbeck, \emph{Quantum Transport through Complex
            Networks}, PhD thesis, Universit\"at Regensburg (2012).
      \item[{[3]}] D.~M.~Wilkins, \emph{A Theoretical Investigation into
            Energy Transfer in Photosynthetic Open Quantum Systems},
            University of Oxford (2012); \texttt{arXiv:1503.03277}.
      \item[{[4]}] A.~Seneviratne, P.~L.~Walters, F.~Wang, \emph{Exact
            Non-Markovian Quantum Dynamics on the NISQ Device Using Kraus
            Operators}, ACS Omega \textbf{9}, 9666 (2024).
      \item[{[5]}] X.~Dan, E.~Geva, V.~S.~Batista, \emph{Simulating
            Non-Markovian Quantum Dynamics on NISQ Computers Using the
            Hierarchical Equations of Motion}, J.~Chem.~Theory Comput.
            \textbf{21}, 1530 (2025).
    \end{itemize}
  \end{block}%
\end{postercolumn}
\end{frame}
\end{document}
